% Generated draft: requires mathematical and source review.
\documentclass[11pt]{article}
\usepackage[margin=25mm]{geometry}
\usepackage{fontspec}
\setmainfont{Latin Modern Roman}
\usepackage{amsmath,amssymb,mathtools,unicode-math}
\setmathfont{Latin Modern Math}
\usepackage{enumitem,needspace,xurl,hyperref}
\setlength{\parindent}{0pt}
\setlength{\parskip}{5pt}
\setlength{\emergencystretch}{3em}
\newcommand{\N}{\mathbb{N}}
\newcommand{\Z}{\mathbb{Z}}
\newcommand{\Q}{\mathbb{Q}}
\newcommand{\R}{\mathbb{R}}
\newcommand{\C}{\mathbb{C}}
\newcommand{\F}{\mathbb{F}}
\newcommand{\A}{\mathbb{A}}
\newcommand{\PP}{\mathbb{P}}
\newcommand{\E}{\mathbb{E}}
\newcommand{\eps}{\varepsilon}
\newcommand{\dd}{\,\mathrm{d}}
\begin{document}
\section*{UnsolvedMath ID: 2301003}
\textbf{Source problem code:} catalogue entry
\subsection{Definitions and mathematical statement}
Let \(f\) be a nonconstant meromorphic function on \(\C\). Write \(n(r,a)\) for the number of \(a\)-points in \(|z|\le r\), counted with multiplicity, using poles for \(a=\infty\), and \(\overline n(r,a)\) for the corresponding count without multiplicity. Let \(N(r,a)\) be the integrated counting function and
\[
T(r,f)=\frac1{2\pi}\int_0^{2\pi}\log^+|f(re^{it})|\,\dd t+N(r,\infty),
\quad
\delta(a,f)=1-\limsup_{r\to\infty}\frac{N(r,a)}{T(r,f)}.
\]
The order is \(\rho=\limsup_{r\to\infty}\log T(r,f)/\log r\); finite order means \(\rho<\infty\). The deficiency sum runs over all \(a\in\widehat\C=\C\cup\{\infty\}\), with nonnegative sums interpreted as suprema of finite subsums.

The conjecture recorded as Problem 1.3 is
\[
\rho<\infty,\quad \sum_{a\in\widehat\C}\delta(a,f)=2
\quad\Longrightarrow\quad
\rho=\frac n2\text{ for an integer }n\ge2,
\qquad \delta(a,f)\in\Q\text{ for every }a.
\]
The source credits F. Nevanlinna with this conclusion under the extra hypothesis of no multiple values, stated there as \(n(r,a)=\overline n(r,a)\) for every value \(a\), and also refers to R. Nevanlinna.

The supplied research summary reports that Drasin (1987) resolves the conjecture, even assuming only finite lower order \(\liminf_{r\to\infty}\log T(r,f)/\log r<\infty\). It also reports that every deficient value is asymptotic, that the reduced denominators of the rational deficiencies are at most \(2\rho\), and a growth representation \(T(r,f)=r^\rho\ell(r)\). Here an asymptotic value means a limit of \(f\) along a curve tending to infinity, in the spherical topology for \(\infty\). The excerpt does not define the auxiliary function \(\ell\), so no additional regularity for it is asserted. The summary mentions a shorter proof and two further properties due to Eremenko without specifying those properties. These are supplied literature claims, not independent verification.
\subsection{Short English statement}
A finite-order meromorphic function with maximal deficiency sum two should have half-integer order at least one and rational deficiencies. The supplied summary reports Drasin's affirmative solution and stronger conclusions.
\subsection{Sources}
\begin{itemize}
\item Local UnsolvedMath registry, permanent ID \(2301003\), exact source code \texttt{catalogue entry}, ``Research Problems in Function Theory --- Problem 1.3.''
\item Hayman, \emph{Research Problems in Function Theory} (2018), Problem 1.3; \url{https://arxiv.org/abs/1809.07200}. Source-TeX extraction; recorded access 2026-07-29. The source attributes an earlier special case to F. Nevanlinna and also mentions R. Nevanlinna, without complete citations here.
\item D. Drasin, ``Proof of a conjecture of F. Nevanlinna concerning functions which have deficiency sum two,'' \emph{Acta Mathematica} \textbf{158} (1987), 1--94, DOI \texttt{10.1007/BF02392256}, as supplied in the research summary.
\item The summary references Eremenko's ``Meromorphic functions of finite order with maximal deficiency sum'' through a MaRDI entry, without a complete citation or URL.
\end{itemize}
\end{document}
